\chapter{Analisi}

\section{Introduzione}

\subsection{Obiettivo}
Si vuole realizzare un’applicazione di messaggistica sicura. L’applicazione
dovra’ gestire un server che fa da tramite e si occupa della ricezione e invio dei messaggi agli
utenti che vogliono comunicare. In questo documento spieghero' nel dettaglio il design,
l'implementazione e le motivazioni delle scelte fatte riguardo le tecnologie utilizzate e gli
algoritmi implementati.

\subsection{Interazioni}
Ci sono 2 tipi di entita', client e server. Ogni client potra' inviare diverse tipologie di messaggi
al server che indicheranno le azioni da fare, ad esempio la registrazione di un nuovo utente,
l'invio di un messaggio, lo scambio di chiavi e diverse altre che vedremo nel dettaglio piu' avanti.

\section{Requisiti}
L'applicazione dovra' garantire:
\begin{itemize}
    \item Confidenzialita': nessuno esterno alla comunicazione deve essere in grado di leggere il
    contenuto dei messaggi scambiati.
    \item Integrita': vogliamo garantire che tutti i messaggi arrivino correttamente.
    \item Identificazione: il server deve essere sicuro dell'identita' dell'utente e viceversa.
\end{itemize}

\section{Tipologie di messaggi}
Ci sono diverse tipologie di messaggi che l'utente puo' mandare al server e vice versa, in questa
sezione tratteremo i tipi di messaggi che l'utente puo' vedere e decidere di inviare, i messaggi che
sono under the hood saranno spiegati nel dettaglio del design e dell'implementazione.

\begin{itemize}
    \item Register [id, password]: il messaggio di register permette di registrare un nuovo utente
    con un id univoco e una password. A questo messaggio corrisponde una risposta del server con
    conferma di esecuzione oppure con rifiuto piu' tipo di errore.
    \item Login [id, password]: il messaggio di login permette di accedere a un account dato l'id e
    la password, anche in questo caso il server rispondera' con un messaggio di successo oppure di
    fallimento con relativo errore.
    \item Send message [id, message]: permette di inviare un messaggio all'utente con l'id
    specificato.
\end{itemize}
I messaggi che il server invia ai client sono tutte risposte a comandi ricevuti oppure servono per
gli scambi di chiavi che saranno approfonditi successivamente.